\chapter{CƠ SỞ LÝ THUYẾT VÀ CÔNG NGHỆ SỬ DỤNG}
\section{Cơ sở lý thuyết về thị giác máy tính và Nhận diện khuôn mặt}

\subsection{Bài toán phát hiện khuôn mặt (Face Detection)}

\subsubsection{Tổng quan về bản chất ma trận điểm ảnh}
Trong thị giác máy tính, một hình ảnh kỹ thuật số không phải là một thực thể trực quan liên tục mà là một tập hợp cấu trúc của các giá trị số rời rạc, được biểu diễn dưới dạng một tensor toán học bậc ba có kích thước $M \times N \times 3$ trong không gian màu RGB (Red - Green - Blue). Mỗi hình ảnh đầu vào bao gồm ba kênh màu độc lập, tương ứng với ba ma trận hai chiều đại diện cho cường độ của các bước sóng ánh sáng đỏ, xanh lá, và xanh dương. Các phần tử $I(x, y, c)$ trong tensor đại diện cho giá trị cường độ sáng của điểm ảnh (pixel) tại tọa độ $(x, y)$ trên kênh màu $c$, thường được mã hóa ở độ phân giải 8-bit với dải giá trị nằm trong đoạn $[0, 255]$.

Khi mô hình mạng nơ-ron tích chập (CNN) xử lý dữ liệu thô này, việc giữ nguyên cấu trúc ảnh màu RGB mang lại những lợi điểm vượt trội so với việc chuyển đổi sang ảnh xám (Grayscale).

Đầu tiên, việc \textbf{bảo toàn thông tin đặc trưng hình học và sắc độ} là rất quan trọng. Không gian màu RGB cung cấp thông tin toàn diện về phân bố màu sắc, sắc độ da, và độ tương phản giữa các vùng ánh sáng. Đây là nguồn dữ liệu quan trọng giúp các kiến trúc mạng TinyML (như MTMN, MobileFaceNet) trích xuất chính xác các đặc trưng cấp thấp (cạnh, góc) và các đặc trưng cấp cao cấu trúc khuôn mặt.

Thứ hai, không gian màu RGB \textbf{hỗ trợ bài toán chống giả mạo khuôn mặt (Anti-spoofing)}. Thông tin màu sắc từ 3 kênh màu giúp mô hình phân biệt được sự khác biệt về chất liệu phản xạ ánh sáng giữa da người thật và các ảnh chụp/màn hình giả mạo, điều mà ảnh đơn kênh Grayscale hoàn toàn bất lực.

Cuối cùng, việc xử lý ảnh màu được \textbf{tối ưu hóa nhờ phần cứng thế hệ mới}. Mặc dù xử lý tensor 3 kênh tốn tài nguyên hơn, nhưng với kiến trúc lõi kép Xtensa LX7 của ESP32-S3 tích hợp bộ tăng tốc vector (AI Vector Extension), các phép toán nhân chập ma trận giữa các bộ lọc (kernels) và tensor RGB đầu vào được thực thi song song bằng phần cứng trực tiếp trên bộ nhớ SRAM. Điều này giúp hệ thống đạt được độ chính xác tối đa mà vẫn đảm bảo thời gian suy luận (inference time) cực thấp mà không cần phải đánh đổi bằng việc hạ cấp không gian màu.

\subsubsection{Nguyên lý hoạt động phân tầng của kiến trúc MTMN}
Mạng tích chập đa nhiệm (Multi-task Cascaded Convolutional Networks - MTCNN) hay biến thể tối ưu hóa cho phần cứng nhúng MTMN là một kiến trúc phân tầng bao gồm ba giai đoạn mạng thần kinh độc lập, được thiết kế chuyên biệt để phát hiện khuôn mặt và các điểm đặc trưng theo cơ chế từ thô đến tinh. Kiến trúc này giải quyết bài toán phát hiện vật thể bằng cách chia nhỏ quá trình sàng lọc ứng viên thông qua ba mạng: P-Net (Proposal Network), R-Net (Refinement Network), và O-Net (Output Network).

Giai đoạn đầu tiên là \textbf{P-Net (Proposal Network)}, sử dụng một mạng tích chập hoàn toàn (Fully Convolutional Network) để tạo ra các vùng ứng viên chứa khuôn mặt tiềm năng. Ảnh đầu vào trước tiên được đưa qua một hệ thống phân cấp tỷ lệ ảnh (Image Pyramid) để tạo ra các bản sao ở nhiều kích thước khác nhau, giúp mạng bắt được các khuôn mặt có quy mô lớn nhỏ khác nhau. P-Net trượt một cửa sổ kích thước $12 \times 12$ trên các cấu trúc ảnh này để thu về các vector đặc trưng, sau đó thực hiện đồng thời hai nhiệm vụ: phân loại thực thể (có phải khuôn mặt hay không) và hồi quy bounding box để tinh chỉnh tọa độ sơ bộ. Kỹ thuật triệt tiêu không tối đa (Non-Maximum Suppression - NMS) được áp dụng ngay sau đó để loại bỏ các vùng chồng lấn có độ tin cậy thấp.

Tiếp theo, các vùng ứng viên vượt qua bộ lọc của P-Net sẽ được chuẩn hóa về kích thước $24 \times 24$ trước khi đưa vào \textbf{R-Net (Refinement Network)}. Vì là một mạng CNN có chứa các lớp kết nối đầy đủ (Fully Connected Layers), R-Net sở hữu năng lực phân biệt mạnh mẽ hơn. Nó thực hiện nhiệm vụ bác bỏ các vùng nhiễu (false positives) do P-Net bỏ sót, tiếp tục hồi quy chính xác hơn tọa độ của bounding box và áp dụng NMS với ngưỡng nghiêm ngặt hơn để giữ lại các ứng viên chất lượng cao.

Giai đoạn cuối cùng là \textbf{O-Net (Output Network)}, nhận đầu vào là các vùng ứng viên đã qua tinh chỉnh và chuẩn hóa chúng về kích thước $48 \times 48$. O-Net có cấu trúc sâu hơn và phức tạp hơn, thực hiện đồng bộ ba nhiệm vụ đa mục tiêu (Multi-task Learning): phân loại thực thể với độ chính xác cao nhất, xác định tọa độ hình học cuối cùng của bounding box, và định vị chính xác tọa độ của 5 điểm neo đặc trưng trên khuôn mặt.

\subsubsection{Ý nghĩa của 5 điểm neo trong việc thực hiện Face Alignment}
Định vị điểm đặc trưng (Facial Landmark Detection) là một bước đệm có tính chất quyết định trước khi đưa dữ liệu vào mạng nhận diện đặc trưng sâu. O-Net xuất ra tọa độ của 5 điểm neo cốt lõi bao gồm: tâm hai mắt, đỉnh mũi, và hai khóe miệng. Ý nghĩa toán học và kỹ thuật của 5 điểm này nằm ở việc thiết lập một phép biến đổi hình học Affine (Affine Transformation) để thực hiện quá trình căn chỉnh khuôn mặt (Face Alignment).

Trong thực tế, khuôn mặt người dùng khi tiếp cận camera của khóa thông minh thường bị nghiêng theo các góc Pitch (gật lên xuống), Yaw (quay trái phải), và Roll (nghiêng hai bên). Sự biến động về góc nhìn này làm thay đổi cấu trúc ma trận điểm ảnh, khiến mạng trích xuất đặc trưng sâu (như MobileFaceNet) dễ bị sai lệch khi so khớp vector. Bằng cách sử dụng tọa độ của 5 điểm neo thu được, hệ thống sẽ tính toán một ma trận biến đổi Affine để dịch chuyển, xoay và thu phóng khuôn mặt thực tế về một hệ tọa độ chuẩn (Template). Quá trình này ép buộc hai mắt phải nằm trên một đường thẳng nằm ngang song song với trục hoành và khoảng cách giữa các thành phần ngũ quan tuân theo một tỷ lệ vàng cố định. Việc loại bỏ hoàn toàn các góc nghiêng Roll và giảm thiểu sai số do góc Yaw gây ra giúp đảm bảo rằng phần ảnh khuôn mặt cắt ra luôn ở góc chính diện, từ đó nâng cao độ chính xác tuyến tính và tính ổn định cho toàn bộ hệ thống nhận diện tại biên.

\subsection{Bài toán trích xuất đặc trưng và nhận diện (Face Recognition)}

\subsubsection{Bản chất toán học của Embedding Vector}
Sau khi khuôn mặt được phát hiện và căn chỉnh (Face Alignment), dữ liệu điểm ảnh thô sẽ được đưa qua một mạng nơ-ron tích chập (CNN) đóng vai trò là một bộ trích xuất đặc trưng (Feature Extractor), điển hình như kiến trúc MobileFaceNet. Về mặt toán học, quá trình này thực chất là một phép ánh xạ phi tuyến tính từ không gian ma trận điểm ảnh nhiều chiều sang một không gian vector liên tục (Latent Space). Phép ánh xạ được định nghĩa là $f: \mathbb{R}^{M \times N \times 3} \rightarrow \mathbb{R}^d$, trong đó đầu ra là một vector đặc trưng (Embedding Vector) có số chiều là $d$.

Trong đề tài này, không gian nhúng được lựa chọn có số chiều là $d = 512$. Sự lựa chọn $512$ chiều không phải là ngẫu nhiên mà là sự cân bằng tối ưu (trade-off) giữa khả năng biểu diễn thông tin và giới hạn tài nguyên của vi điều khiển. Nếu sử dụng số chiều nhỏ (như $64$ hay $128$), không gian toán học sẽ không đủ độ phức tạp để phân tách ranh giới tuyến tính giữa hàng triệu đặc điểm khuôn mặt khác nhau, dẫn đến hiện tượng trùng lặp (collision) và làm giảm độ chính xác (Accuracy). Ngược lại, nếu sử dụng $1024$ hoặc $2048$ chiều, khối lượng phép tính dấu phẩy động (FLOPs) và dung lượng RAM cần thiết để lưu trữ sẽ vượt quá khả năng của bộ nhớ SRAM 512 KB trên hệ thống ESP32-S3.

Về mặt ý nghĩa vật lý, 512 chiều này không biểu diễn trực tiếp các thông số sinh trắc học có thể đo đếm bằng mắt thường (như khoảng cách giữa hai mắt, độ rộng của mũi). Thay vào đó, chúng biểu diễn các mức độ kích hoạt của các bộ lọc trừu tượng đối với các cấu trúc vi mô (micro-structures), kết cấu bề mặt da (texture), và độ phản xạ ánh sáng (albedo) tại các khu vực khác nhau trên khuôn mặt. Mỗi khuôn mặt sau khi đi qua mạng sẽ được định vị thành một điểm duy nhất trong không gian hình học 512 chiều này.

\paragraph{Công thức toán học của các thang đo khoảng cách}
Để xác định hai khuôn mặt có thuộc về cùng một người hay không, hệ thống cần tính toán độ tương đồng giữa hai Embedding Vector $\mathbf{A}$ và $\mathbf{B}$ thu được. Hai phương pháp đo lường toán học phổ biến nhất trong không gian vector là Khoảng cách Euclid (Euclidean Distance) và Độ tương đồng Cosine (Cosine Similarity).

Khoảng cách Euclid đo lường khoảng cách hình học tuyệt đối theo đường thẳng giữa hai điểm trong không gian $\mathbb{R}^d$, được định nghĩa bằng công thức chuẩn L2 (L2-Norm):
$$
d(\mathbf{A}, \mathbf{B}) = \|\mathbf{A} - \mathbf{B}\|_2 = \sqrt{\sum_{i=1}^{d} (A_i - B_i)^2}
$$
Trong đó, $A_i$ và $B_i$ là giá trị tại chiều thứ $i$ của hai vector. Khoảng cách càng nhỏ (tiến về 0) thì hai khuôn mặt càng giống nhau.

Độ tương đồng Cosine đo lường góc $\theta$ được tạo bởi hai vector trong không gian đa chiều, không quan tâm đến độ lớn (magnitude) của vector mà chỉ xét hướng (direction) của chúng. Công thức được tính bằng tỷ số giữa tích vô hướng và tích của độ dài hai vector:
$$
S_C(\mathbf{A}, \mathbf{B}) = \cos(\theta) = \frac{\mathbf{A} \cdot \mathbf{B}}{\|\mathbf{A}\| \|\mathbf{B}\|} = \frac{\sum_{i=1}^{d} A_i B_i}{\sqrt{\sum_{i=1}^{d} A_i^2} \sqrt{\sum_{i=1}^{d} B_i^2}}
$$
Giá trị của Cosine Similarity nằm trong đoạn $[-1, 1]$. Khi hai khuôn mặt giống hệt nhau, góc giữa chúng bằng $0^\circ \Rightarrow \cos(0) = 1$.

\paragraph{Phân tích và lựa chọn giải pháp tối ưu cho vi điều khiển}
Đối với một hệ thống máy chủ có vi xử lý đồ họa (GPU) mạnh mẽ, việc lựa chọn giữa Euclid hay Cosine không gây ra sự chênh lệch lớn về mặt thời gian thực thi. Tuy nhiên, khi triển khai mô hình lên các thiết bị tính toán tại biên (Edge Device) như vi điều khiển ESP32-S3, yếu tố tài nguyên phần cứng mang tính quyết định. Đề tài quyết định lựa chọn \textbf{Độ tương đồng Cosine (Cosine Similarity)} kết hợp với chuẩn hóa vector làm phương pháp so khớp cốt lõi dựa trên các luận điểm sau:

Thứ nhất, phép tính khoảng cách Euclid yêu cầu vi điều khiển phải thực hiện phép trừ, phép bình phương, tổng quát hóa, và nặng nề nhất là phép khai căn bậc hai ($\sqrt{x}$). Đối với kiến trúc phần cứng của vi điều khiển, việc giải quyết các bài toán khai căn bậc hai cho số thực dấu phẩy động (Floating-point) tiêu tốn rất nhiều chu kỳ xung nhịp (clock cycles), gây ra độ trễ (latency) lớn trong quá trình nhận diện.

Thứ hai, khi sử dụng Cosine Similarity, hệ thống được thiết kế để áp dụng phép chuẩn hóa L2 (L2-Normalization) ngay tại đầu ra của mạng CNN. Điều này ép độ dài của mọi vector về $1$ (tức là $\|\mathbf{A}\| = 1$ và $\|\mathbf{B}\| = 1$). Khi đó, mẫu số trong công thức Cosine sẽ bằng $1$, và phép toán phức tạp ban đầu được triệt tiêu hoàn toàn, chỉ còn lại một phép nhân tích vô hướng (Dot Product) duy nhất:
$$
S_C(\mathbf{A}, \mathbf{B}) \approx \sum_{i=1}^{d} A_i B_i
$$
Hơn thế nữa, lõi vi xử lý Xtensa LX7 trên ESP32-S3 được trang bị bộ tập lệnh mở rộng (AI/DSP instructions) tối ưu hóa đặc biệt cho các phép toán Multiply-Accumulate (MAC). Vi điều khiển có thể xử lý phép tính $\sum A_i B_i$ chỉ trong vài chu kỳ máy ở cấp độ thanh ghi (Register-level) mà không cần phải truy xuất liên tục vào bộ nhớ SRAM. Sự tối ưu hóa về mặt thuật toán toán học kết hợp với kiến trúc phần cứng này giúp hệ thống đạt được tốc độ nhận diện thời gian thực (Real-time Face Recognition) ngay cả với cơ sở dữ liệu lên tới 100 người dùng lưu trữ cục bộ.

\subsection{Tổng quan về TinyML (Học máy trên thiết bị nhúng)}

\subsubsection{Phân tích các thách thức phần cứng (Hardware Bottlenecks)}
Trí tuệ nhân tạo vạn vật (Artificial Intelligence of Things - AIoT) đang chuyển dịch mạnh mẽ từ điện toán đám mây xuống các thiết bị tại biên thông qua khái niệm TinyML (Học máy quy mô nhỏ). Tuy nhiên, việc triển khai các mạng nơ-ron nhân tạo lên vi điều khiển phải đối mặt với những nút thắt cổ chai (bottlenecks) cực kỳ nghiêm trọng về kiến trúc phần cứng, đòi hỏi sự tối ưu hóa khắt khe từ cấp độ bộ nhớ đến kiến trúc tập lệnh.

Thách thức đầu tiên và lớn nhất là giới hạn về dung lượng và tốc độ của bộ nhớ. Vi điều khiển xử lý học sâu cần nạp trọng số (weights) và các ma trận kích hoạt (activations) vào bộ nhớ RAM tĩnh (Static-RAM) để tính toán. Trên ESP32-S3, dung lượng bộ nhớ SRAM nội bộ chỉ đạt 512 KB, một con số quá nhỏ bé so với hàng chục Megabyte của các mô hình CNN thông thường. Để giải quyết, các nhà thiết kế hệ thống thường mở rộng bằng RAM giả tĩnh (Pseudo-Static RAM) bên ngoài với dung lượng lớn hơn, có thể lên tới 8 MB hoặc 16 MB. Tuy nhiên, đặc tính vật lý của PSRAM lại tạo ra một điểm nghẽn mới: tốc độ truy xuất. PSRAM giao tiếp với nhân vi xử lý thông qua bus ngoại vi (thường là Quad SPI hoặc Octal SPI), dẫn đến độ trễ truy cập cao hơn rất nhiều so với SRAM nội bộ. Khi hàng triệu phép nhân ma trận liên tục gọi dữ liệu từ PSRAM, CPU sẽ rơi vào trạng thái chờ, làm giảm nghiêm trọng tốc độ khung hình (FPS) của luồng nhận diện. Ngoài ra, việc lưu trữ toàn bộ kiến trúc mạng trước khi nạp vào RAM cũng bị giới hạn bởi bộ nhớ Flash tĩnh, yêu cầu mô hình phải cực kỳ nhỏ gọn để chừa không gian cho hệ điều hành thời gian thực (RTOS) và các phân vùng hệ thống khác (NVS, SPIFFS).

Thách thức thứ hai đến từ năng lực tính toán và vấn đề năng lượng. Không giống như các máy chủ được trang bị vi xử lý đồ họa (GPU) chuyên dụng cho xử lý song song, vi điều khiển hoàn toàn dựa vào bộ xử lý trung tâm (CPU) để thực hiện các phép toán học sâu. Việc thiếu hụt lõi tính toán Tensor độc lập khiến vi điều khiển phải gánh chịu tải tính toán cực lớn. Kéo theo đó là sự gia tăng đột biến về nhiệt năng tỏa ra. Trong các ứng dụng thực tế như khóa cửa thông minh, bo mạch thường được đóng gói kín trong các không gian hẹp không có hệ thống tản nhiệt chủ động. Nhiệt độ gia tăng làm giảm hiệu suất của cảm biến camera (tạo ra nhiễu nhiệt trên điểm ảnh) và có thể kích hoạt cơ chế tự giảm xung nhịp của vi điều khiển, làm suy giảm tốc độ nhận diện hoặc thậm chí gây sụt áp, treo hệ thống.

\subsubsection{Nguyên lý của kỹ thuật Lượng tử hóa (Quantization)}
Để vượt qua rào cản về bộ nhớ và năng lực tính toán, TinyML áp dụng triệt để kỹ thuật Lượng tử hóa (Quantization). Đây là quá trình chuyển đổi biểu diễn dữ liệu của các trọng số mạng và giá trị kích hoạt từ định dạng số thực dấu phẩy động 32-bit (Float32) sang định dạng số nguyên 8-bit (INT8).

Việc chuyển đổi từ Float32 (chiếm 4 byte) sang INT8 (chiếm 1 byte) mang lại lợi ích kép: giảm trực tiếp 4 lần dung lượng bộ nhớ lưu trữ (cả Flash và RAM) và cho phép các phép toán nhân cộng tích lũy (MAC) thực hiện trực tiếp trên bộ số học và logic nguyên thay vì bộ xử lý dấu phẩy động, qua đó tiết kiệm đáng kể chu kỳ máy và năng lượng.

Về mặt toán học, lượng tử hóa đồng đều (Uniform Quantization) thực hiện phép ánh xạ một dải giá trị thực liên tục $r$ sang một dải giá trị nguyên rời rạc $q$ (thường nằm trong đoạn $[-128, 127]$ đối với INT8 có dấu) thông qua công thức ánh xạ tuyến tính:
$$
r = S(q - Z)
$$
Trong đó $r$ là giá trị thực (Float32) ban đầu của trọng số hoặc đầu ra của lớp kích hoạt, $q$ là giá trị số nguyên sau khi được lượng tử hóa (INT8), $S$ (Scale) là hệ số co giãn, một giá trị Float32 dương duy nhất đại diện cho khoảng cách giữa hai bước lượng tử kề nhau và $Z$ (Zero-point) là điểm tham chiếu số nguyên tương ứng với giá trị thực $r = 0$, giúp bảo toàn tính chất của các lớp đệm (padding) và các hàm kích hoạt (như ReLU) luôn lấy số 0 làm bản lề.

Từ công thức trên, quá trình chuyển đổi ngược từ số thực sang số nguyên (để lưu trữ vào chip) được tính toán như sau (kết hợp với hàm làm tròn số):
$$
q = \text{round}\left(\frac{r}{S} + Z\right)
$$
Mặc dù việc ép hàng triệu giá trị thập phân vào 256 mức rời rạc của INT8 gây ra sai số lượng tử hóa (Quantization Error), các kỹ thuật lượng tử hóa hiện đại (như Post-Training Quantization) có khả năng hiệu chuẩn các tham số $S$ và $Z$ thông qua một tập dữ liệu đại diện nhỏ. Quá trình này giúp mô hình khôi phục và bảo toàn gần như nguyên vẹn độ chính xác (Accuracy) so với mô hình Float32 gốc, biến TinyML trở thành giải pháp cốt lõi để đưa các thuật toán AI phức tạp lên thiết bị nhúng một cách hiệu quả và khả thi nhất.

\section{Công nghệ phần cứng sử dụng}

Trong phạm vi hệ thống này, việc lựa chọn kiến trúc phần cứng đóng vai trò quyết định đến hiệu năng xử lý thuật toán tại biên (Edge Computing) và độ tin cậy của hệ thống khóa thông minh. Phân mục này sẽ phân tích chuyên sâu về vi điều khiển lõi trung tâm ESP32-S3 CAM cùng mô-đun cảm biến chạm ngoại vi.

\subsection{Mô-đun điều khiển trung tâm: ESP32-S3 CAM}

Mô-đun ESP32-S3 CAM là sự kết hợp giữa vi xử lý SoC ESP32-S3 mạnh mẽ của Espressif Systems và giao tiếp ngoại vi camera, được tối ưu hóa đặc biệt cho các tác vụ trí tuệ nhân tạo thị giác máy tính ở cấp độ nhúng (TinyML).

\subsubsection{Kiến trúc lõi kép Xtensa LX7 và Tập lệnh mở rộng Vector (Vector Extension):}
ESP32-S3 được trang bị bộ vi xử lý lõi kép Xtensa® 32-bit LX7, vận hành với xung nhịp lên tới 240 MHz. Kiến trúc đường ống (pipeline) 7 giai đoạn của LX7 cho phép tối ưu hóa tốc độ thực thi lệnh trên mỗi chu kỳ xung nhịp. Tuy nhiên, điểm cải tiến mang tính đột phá của ESP32-S3 so với các thế hệ tiền nhiệm (như ESP32 lõi LX6) chính là việc tích hợp sẵn tập lệnh mở rộng Vector (Vector Extension) và bộ tăng tốc phần cứng cho các phép toán ma trận.

Trong các mô hình mạng thần kinh nhân tạo sâu (CNN) phục vụ bài toán nhận diện khuôn mặt, các phép toán cốt lõi chiếm tới hơn 90\% thời gian tính toán là phép nhân và cộng tích lũy (Multiply-Accumulate - MAC) trên các ma trận trọng số (weights) và ma trận đặc trưng (feature maps). Tập lệnh Vector trên ESP32-S3 hỗ trợ kiến trúc SIMD (Single Instruction, Multiple Data), cho phép một câu lệnh duy nhất xử lý đồng thời nhiều dữ liệu định dạng số nguyên 8-bit hoặc 16-bit định điểm tĩnh (fixed-point). Bộ tăng tốc này giúp tối ưu hóa đáng kể các hàm kích hoạt (Activation Functions) và các lớp tích chập (Convolutional Layers). Kết quả là, khi kết hợp với thư viện tối ưu hóa phần cứng ESP-DL, tốc độ tính toán các phép nhân ma trận tăng từ 4 đến 10 lần so với việc thực thi bằng phần mềm thuần túy trên CPU lõi thông thường, đáp ứng thời gian thực (real-time inference) cho luồng dữ liệu video từ camera.

\subsubsection{Phân tích cấu trúc bộ nhớ SRAM và PSRAM trong tác vụ xử lý ảnh:}
Xử lý dữ liệu thị giác máy tính yêu cầu một kiến trúc bộ nhớ phân cấp luân chuyển băng thông lớn. ESP32-S3 sở hữu 512 KB bộ nhớ SRAM nội chip (Internal SRAM). Bộ nhớ SRAM này vận hành ở tốc độ bằng với xung nhịp CPU, độ trễ bằng không (zero wait-state), cực kỳ thích hợp để lưu trữ mã lệnh thực thi nhanh, các biến cục bộ và các mảng dữ liệu tính toán tức thời của thuật toán AI. Tuy nhiên, dung lượng 512 KB là hoàn toàn bất khả thi để chứa toàn bộ một khung ảnh độ phân giải cao cùng các tham số mô hình nhận diện.

Để giải quyết bài toán thiếu hụt bộ nhớ, mô-đun ESP32-S3 CAM tích hợp thêm bộ nhớ PSRAM ngoại vi (Pseudo-SRAM) mở rộng (thường là 2 MB hoặc 8 MB tùy phiên bản) kết nối qua giao tiếp SPI cấu hình bát phân (Octal SPI) tốc độ cao. Sự phân biệt vai trò giữa hai loại bộ nhớ này trong tác vụ xử lý ảnh từ Camera được mô tả cụ thể. \textbf{SRAM nội chip} được sử dụng làm nơi triển khai các lớp tính toán trung gian của mạng CNN, lưu trữ các vector đặc trưng khuôn mặt (embedding vectors) sau khi trích xuất (thường có kích thước nhỏ, ví dụ 512 float32). Do tốc độ truy xuất cực cao, SRAM đảm bảo các vòng lặp toán học của thuật toán nhận diện không bị nghẽn cổ chai (bottleneck) về băng thông dữ liệu. Trong khi đó, \textbf{PSRAM ngoại vi} đóng vai trò là phân vùng lưu trữ cho bộ đệm khung hình (Frame Buffer) của Camera. Một khung ảnh thô định dạng RGB565 độ phân giải QVGA ($320 \times 240$) cần khoảng 153.6 KB, nhưng nếu cấu hình lên độ phân giải VGA ($640 \times 480$), dung lượng yêu cầu lên tới 614.4 KB, vượt quá giới hạn tổng của SRAM nội chip. Hơn nữa, quá trình giải mã ảnh JPEG hoặc lưu trữ các ma trận trọng số của mô hình TinyML đòi hỏi không gian nhớ lên tới vài Megabyte. Cơ chế ánh sáng Cache (Cache Mapping MMU) của ESP32-S3 cho phép CPU truy cập vào PSRAM thông qua không gian địa chỉ ảo một cách mượt mà, giúp hệ thống xử lý được các tác vụ xử lý ảnh phức tạp mà không gặp lỗi tràn bộ nhớ (Out of Memory).

\subsubsection{Cảm biến ảnh SoC độ phân giải cao OV3660}

Trong hệ thống này, cảm biến ảnh OV3660 của OmniVision được lựa chọn làm phần cứng thu nhận tín hiệu thị giác chính thức. Đây là dòng cảm biến ảnh CMOS SoC (System on Chip) sở hữu cấu trúc phần cứng mạnh mẽ và bộ xử lý tín hiệu tích hợp sâu, cung cấp dữ liệu hình ảnh chất lượng cao cho các thuật toán nhận diện khuôn mặt tại biên.

\paragraph{Thông số quang học và công nghệ kiến trúc điểm ảnh:}
OV3660 hỗ trợ độ phân giải tối đa lên tới 3 Megapixel ($2048 \times 1536 \text{ pixels}$). Việc sở hữu mật độ điểm ảnh lớn giúp hệ thống thu thập được các trường đặc trưng hình học sắc nét và chi tiết trên khuôn mặt đối tượng ngay cả ở khoảng cách xa. Cảm biến ứng dụng công nghệ kiến trúc điểm ảnh OmniBSI™ (Backside Illumination), dịch chuyển toàn bộ các đường mạch kim loại dẫn tín hiệu ra phía sau lớp cảm quang. Thiết kế này tối ưu hóa hiệu suất lượng tử (Quantum Efficiency), tăng tối đa khả năng hấp thụ ánh sáng trên mỗi micron vuông của điểm ảnh, giúp hệ thống hoạt động ổn định trong các môi trường thiếu sáng hoặc có luồng sáng phức tạp của không gian thực tế.

Dữ liệu thô sau khi qua mảng lọc màu Bayer sẽ được xử lý và xuất ra dưới nhiều định dạng linh hoạt như RAW RGB, RGB565, YUV422/420. Đặc biệt, chip tích hợp sẵn bộ nén phần cứng JPEG chất lượng cao, giúp nén luồng dữ liệu hình ảnh trực tiếp trước khi truyền về vi điều khiển, làm giảm áp lực đáng kể lên băng thông truyền tải và dung lượng bộ đệm khung hình (Frame Buffer).

\paragraph{Bộ xử lý tín hiệu hình ảnh (ISP) tích hợp nội chip:}
Ưu thế vượt trội của OV3660 nằm ở bộ xử lý ISP tích hợp ngay trên tầng phần cứng của cảm biến, tự động thực hiện chu trình tiền xử lý ảnh trước khi đổ dữ liệu vào bộ nhớ của ESP32-S3. Các chức năng chính bao gồm \textbf{AEC (Automatic Exposure Control) và AGC (Automatic Gain Control)}, tự động tối ưu hóa thời gian phơi sáng và hệ số khuếch đại analog, ngăn chặn hiện tượng cháy sáng hoặc tối đen. \textbf{AWB (Automatic White Balance)} giúp cân bằng trắng tự động dựa trên nhiệt độ màu môi trường, bảo toàn tính chính xác của sắc độ da. \textbf{LCC (Lens Correction) và DPC (Defective Pixel Correction)} tự động khử hiện tượng tối góc và cô lập các điểm ảnh lỗi. Cuối cùng, \textbf{mạch lọc nhiễu không gian (Advanced Denoising)} loại bỏ nhiễu hạt kỹ thuật số nhưng vẫn giữ lại độ sắc nét của các đường biên cạnh, phục vụ tối ưu cho việc định vị các điểm mốc khuôn mặt.

\paragraph{Nguyên lý giao tiếp DVP (Digital Video Port) với ESP32-S3:}
OV3660 kết nối song song và đồng bộ với bộ điều khiển ngoại vi của ESP32-S3 thông qua chuẩn giao tiếp DVP. Chu trình truyền tải này đòi hỏi sự phối hợp nghiêm ngặt về mặt thời gian giữa các đường tín hiệu. \textbf{Bus dữ liệu song song (D0 - D7)} gồm 8 đường truyền độc lập, chịu trách nhiệm tải giá trị byte dữ liệu của từng pixel. \textbf{Xung nhịp hệ thống (XCLK)} do ESP32-S3 cung cấp để duy trì hoạt động cho cảm biến. \textbf{Xung nhịp điểm ảnh (PCLK)} do OV3660 phát ngược lại, và tại mỗi cạnh lên của PCLK, ESP32-S3 sẽ chốt dữ liệu. \textbf{Tín hiệu đồng bộ dòng (HREF)} định vị phạm vi hợp lệ của một hàng pixel. \textbf{Tín hiệu đồng bộ khung (VSYNC)} đánh dấu điểm bắt đầu và kết thúc của một khung ảnh hoàn chỉnh, kích hoạt ngắt phần cứng trên ESP32-S3 để bộ điều khiển DMA tự động ghi dữ liệu ảnh vào PSRAM. Cuối cùng, \textbf{giao tiếp cấu hình SCCB (SIOD / SIOC)}, tương thích với chuẩn I2C, được dùng để ESP32-S3 truy cập và ghi các giá trị cấu hình vào hệ thống thanh ghi nội bộ của OV3660.

\paragraph{Tối ưu hóa cho bài toán nhận diện và chống giả mạo khuôn mặt:}
Việc sử dụng trực tiếp cảm biến 3MP như OV3660 mang lại lợi thế kỹ thuật rất lớn cho bài toán TinyML. Độ phân giải cao giúp bảo toàn các chi tiết tần số cao trên bề mặt vật thể. Khi kết hợp với các thuật toán phân tích texture da hoặc chiều sâu trường ảnh, hệ thống không chỉ nhận diện danh tính chính xác (giảm tỷ lệ sai lệch FAR/FRR) mà còn hỗ trợ đắc lực cho các module chống giả mạo khuôn mặt (Anti-spoofing), giúp phát hiện và ngăn chặn các hành vi gian lận bằng ảnh in màu hay màn hình kỹ thuật số tái tạo.

\subsection{Linh kiện ngoại vi: Cảm biến chạm (Touch Sensor)}

Trong hệ thống khóa thông minh, cảm biến chạm đóng vai trò làm giao diện tương tác trực tiếp giữa người dùng và hệ thống điều khiển, thay thế cho các nút nhấn cơ học truyền thống dễ bị hao mòn vật lý và ảnh hưởng bởi môi trường ẩm mốc.

\subsubsection{Nguyên lý cảm ứng điện dung (Capacitive Touch Sensing):}
Cảm biến chạm hoạt động dựa trên nguyên lý thay đổi điện dung của một điện cực khi có vật thể dẫn điện (như ngón tay người) tiếp cận. Về mặt vật lý, tấm đồng cảm biến trên bo mạch đóng vai trò như một bản cực của tụ điện, điện môi xung quanh là không khí và lớp vỏ bảo vệ. Điện dung ký sinh ban đầu của hệ thống đối với mặt đất (Ground) được ký hiệu là $C_0$, tuân theo công thức cơ bản:
\begin{equation}
    C = \epsilon \cdot \frac{A}{d}
\end{equation}
Trong đó $\epsilon$ là hằng số điện môi, $A$ là diện tích bề mặt điện cực, và $d$ là khoảng cách giữa các bản cực.

Khi ngón tay người (vốn mang đặc tính dẫn điện và được nối đất tương đối thông qua cơ thể) chạm vào hoặc đến gần vùng cảm biến, một điện dung bổ sung $C_{finger}$ sẽ được hình thành song song với điện dung ký sinh ban đầu. Lúc này, tổng điện dung tương đương tại chân cảm biến tăng lên:
\begin{equation}
    C_{total} = C_0 + C_{finger}
\end{equation}
Sự gia tăng điện dung $\Delta C = C_{finger}$ này sẽ làm thay đổi hằng số thời gian phóng nạp RC của mạch dao động tích hợp bên trong IC cảm biến (ví dụ dòng chip chuyên dụng TTP223 hoặc bộ chuyển đổi cảm ứng nội chip của chính ESP32-S3). Mạch định thời bên trong sẽ đo đạc sự thay đổi tần số dao động hoặc thời gian nạp đầy tụ. Khi mức thay đổi vượt qua một ngưỡng (Threshold) đã được thiết lập sẵn, IC sẽ xác nhận có sự kiện "chạm" (Touch Event).

\subsubsection{Cấu trúc mạch và cơ chế giao tiếp:}
Mô-đun cảm biến chạm thường sử dụng IC TTP223 cấu hình ở chế độ ngõ ra số (Digital Output). Mạch bao gồm các thành phần chính. \textbf{Điện cực cảm ứng (Touch Pad)} thường là một vùng phủ đồng diện tích khoảng $1 \times 1 \text{ cm}^2$ trên bo mạch PCB, được phủ lớp sơn hàn cách điện. \textbf{Tụ tinh chỉnh độ nhạy ($C_s$)}, thường có giá trị từ 0 pF đến 50 pF, được mắc song song từ chân cảm biến xuống Ground để hiệu chuẩn thiết bị cho các lớp vỏ bảo vệ dày. \textbf{Đường tín hiệu điều khiển (VCC, GND, OUT)} kết nối trực tiếp đến một chân GPIO của ESP32-S3. Khi không có tương tác, chân OUT giữ ở mức thấp; khi phát hiện chạm, chân OUT chuyển sang mức cao. ESP32-S3 sử dụng cơ chế ngắt cạnh lên trên chân GPIO này để đánh thức hệ thống từ chế độ ngủ sâu hoặc kích hoạt chu trình nhận diện, giúp tối ưu hóa điện năng.

\section{Công nghệ phần mềm và Thư viện}

\subsection{ESP-IDF và Môi trường phát triển}
Quá trình phát triển phần mềm cho hệ thống nhúng được thực hiện chủ yếu trên nền tảng ESP-IDF kết hợp cùng ngôn ngữ C/C++. Điểm cốt lõi làm nên sức mạnh của hệ sinh thái này là sự tích hợp sâu của Hệ điều hành thời gian thực FreeRTOS. Khác với các vòng lặp tuần tự đơn luồng truyền thống, FreeRTOS cho phép vi điều khiển ESP32-S3 phân chia tài nguyên CPU lõi kép để thực thi đồng thời nhiều tác vụ độc lập thông qua cơ chế lập lịch ưu tiên ưu tiên ngắt.

Trong kiến trúc của hệ thống khóa thông minh, việc quản lý luồng dữ liệu đòi hỏi giải quyết bài toán xung đột tài nguyên giữa hai nhóm tác vụ chính là xử lý thị giác máy tính và duy trì kết nối mạng. Tác vụ thu thập khung hình từ cảm biến camera và chạy suy luận mạng nơ-ron đòi hỏi khối lượng tính toán cực lớn, dễ dàng chiếm dụng toàn bộ chu kỳ máy và gây ra hiện tượng nghẽn cổ chai. Nếu hệ thống hoạt động đơn luồng, quá trình tính toán ma trận kéo dài hàng trăm mili-giây sẽ chặn đứng các ngắt giao tiếp mạng, khiến kết nối Wi-Fi bị ngắt quãng hoặc giao thức MQTT bị rớt do không thể phản hồi gói tin duy trì kết nối.

FreeRTOS giải quyết triệt để vấn đề này bằng cách phân bổ các tác vụ lên hai lõi vật lý của vi xử lý Xtensa LX7. Lõi số không được ưu tiên cấu hình để xử lý ngăn xếp mạng Wi-Fi và duy trì kết nối liên tục với MQTT Broker. Lõi số một được dành riêng cho việc giao tiếp với giao diện DVP của camera, đưa ảnh vào bộ nhớ PSRAM và thực thi các hàm API của thư viện học sâu. Bộ lập lịch của FreeRTOS sẽ tự động tạm dừng các tác vụ xử lý ảnh có độ ưu tiên thấp hơn để nhường quyền điều khiển CPU cho ngắt mạng khi có gói tin đến, đảm bảo hệ thống luôn trong trạng thái trực tuyến sẵn sàng nhận lệnh mở cửa từ xa. Sự đồng bộ dữ liệu giữa các luồng tác vụ này được đảm bảo an toàn thông qua các cơ chế hàng đợi tin nhắn và cờ hiệu, giúp luồng video và tín hiệu cảnh báo được truyền tải mượt mà lên ứng dụng di động mà không gây sụp đổ hệ thống.

\subsection{Khung phần mềm ESP-WHO và thư viện ESP-DL}
Để triển khai thành công các thuật toán trí tuệ nhân tạo lên thiết bị có tài nguyên hạn chế, hệ thống ứng dụng khung phần mềm ESP-WHO chuyên dụng cho xử lý ảnh kết hợp cùng lõi tính toán ESP-DL. ESP-DL đóng vai trò là một thư viện toán học cấp thấp, cung cấp các hàm giao diện lập trình ứng dụng để thực thi các lớp mạng nơ-ron được tối ưu hóa trực tiếp bằng hợp ngữ cho kiến trúc tập lệnh vector trên chip ESP32-S3.

Sức mạnh của ESP-DL nằm ở khả năng tăng tốc phần cứng đối với các phép toán nhân cộng tích lũy cường độ cao. Các lớp toán học cơ bản cấu thành nên mạng học sâu như lớp tích chập, lớp tích chập tách biệt chiều sâu, lớp gộp và lớp kết nối đầy đủ đều được tái cấu trúc bộ nhớ để phù hợp với thanh ghi độ rộng lớn của vi điều khiển. Thay vì xử lý từng điểm dữ liệu đơn lẻ, ESP-DL sử dụng các lệnh xử lý đa dữ liệu đơn luồng lệnh để tính toán đồng thời nhiều phần tử của ma trận trọng số và ma trận đặc trưng. Kỹ thuật này giảm thiểu tối đa số lần truy xuất bộ nhớ ngoài, giúp thuật toán vượt qua rào cản tốc độ truyền tải chậm của băng thông PSRAM, từ đó rút ngắn thời gian suy luận trên mỗi khung hình.

Trên nền tảng tăng tốc toán học đó, hệ thống sử dụng MobileFaceNet làm kiến trúc mạng chủ đạo cho bài toán trích xuất đặc trưng khuôn mặt. MobileFaceNet là một mạng nơ-ron tích chập tinh gọn được thiết kế đặc biệt cho các thiết bị di động và vi điều khiển có giới hạn nghiêm ngặt về dung lượng lưu trữ. Kiến trúc này kế thừa nền tảng khối dư đảo ngược từ MobileNetV2 kết hợp với các cấu trúc nút thắt cổ chai tuyến tính, giúp duy trì độ sâu của mạng trong khi giảm thiểu hàng triệu tham số so với các mô hình truyền thống như ResNet.

Điểm đột phá nhất của MobileFaceNet khi áp dụng vào hệ thống nhúng là việc loại bỏ hoàn toàn lớp kết nối đầy đủ hoặc lớp gộp trung bình toàn cục ở tầng cuối cùng của mạng. Thay vào đó, kiến trúc này sử dụng một lớp tích chập tách biệt chiều sâu toàn cục để tổng hợp các bản đồ đặc trưng không gian thành một vector nhúng đa chiều. Giải pháp này giúp cắt giảm hơn phân nửa tổng số lượng tham số tính toán của toàn bộ mô hình nhưng vẫn giữ được khả năng biểu diễn thông tin sinh trắc học với độ chính xác cao. Việc sử dụng MobileFaceNet chạy trên nền tảng thư viện ESP-DL đảm bảo mô hình có thể nằm gọn trong bộ nhớ tĩnh của thiết bị, đồng thời cung cấp tốc độ trích xuất đặc trưng vượt trội, đáp ứng yêu cầu nhận diện tức thời của thiết bị khóa thông minh.

\subsection{Nền tảng quản lý, giao thức giao tiếp và cơ sở dữ liệu}

\subsubsection{Framework đa nền tảng Flutter}
Flutter là một bộ công cụ phát triển phần mềm nguồn mở (SDK) do Google phát triển, cho phép xây dựng các ứng dụng biên dịch nguyên bản (native compiled) cho cả thiết bị di động (iOS, Android), Web và Desktop từ một kho mã nguồn duy nhất (Single Codebase). Thay vì sử dụng các thành phần giao diện (UI widgets) gốc của hệ điều hành, Flutter sở hữu một engine đồ họa riêng mang tên Skia (hoặc Impeller) để tự vẽ trực tiếp mọi pixel trên màn hình với tốc độ lên đến 60 hoặc 120 khung hình/giây (FPS).

Ngôn ngữ lập trình được sử dụng là Dart, một ngôn ngữ hướng đối tượng mạnh mẽ, hỗ trợ cơ chế biên dịch AOT (Ahead-Of-Time) giúp tối ưu hóa hiệu năng thực thi tương đương với mã nguồn native, kết hợp với tính năng JIT (Just-In-Time) phục vụ cơ chế "Hot Reload", tăng tốc đáng kể quá trình phát triển ứng dụng di động quản lý khóa thông minh.

\subsubsection{Framework Backend NestJS và Prisma ORM}
Tại máy chủ trung tâm (Server), hệ thống sử dụng NestJS – một framework Node.js tiến tiến được xây dựng trên nền tảng TypeScript. NestJS áp dụng kiến trúc Module, Dependency Injection (DI) chặt chẽ lấy cảm hứng từ Angular, giúp mã nguồn backend có tính cấu trúc cao, dễ bảo trì và mở rộng.

Để tương tác với hệ quản trị cơ sở dữ liệu, NestJS được tích hợp kết hợp với Prisma – một công cụ ánh xạ quan hệ - đối tượng (ORM) thế hệ mới. Prisma tự động sinh ra các kiểu dữ liệu (Types) an toàn dựa trên file đặc tả cấu trúc `schema.prisma`, giúp lập trình viên thao tác với cơ sở dữ liệu thông qua các hàm đối tượng thuần túy mà không cần viết các câu lệnh SQL thô, đồng thời tối ưu hóa hiệu năng truy vấn thông qua cơ chế gom cụm (query batching) tự động.

\subsubsection{Mô hình truyền thông kết hợp: REST API, WebSocket và MQTT}
Để đảm bảo tính toàn vẹn thông tin cũng như tốc độ phản hồi tức thời của khóa thông minh, hệ thống triển khai mô hình giao tiếp hỗn hợp tối ưu hóa cho từng loại tác vụ nghiệp vụ.

\textbf{REST API (Representational State Transfer)} hoạt động trên giao thức HTTP/HTTPS, tuân thủ mô hình Request-Response truyền thống. REST API được sử dụng cho các tác vụ quản lý dữ liệu tĩnh hoặc bất đồng bộ không yêu cầu thời gian thực (tác vụ CRUD), chẳng hạn như đăng ký thông tin cư dân, chỉnh sửa quyền truy cập hoặc xem lại lịch sử ra vào lưu trên hệ thống.

\textbf{WebSocket} là giao thức truyền thông cho phép thiết lập kênh giao tiếp song công toàn phần (Full-duplex) liên tục qua một kết nối TCP duy nhất sau cái bắt tay HTTP (Handshake). Trong đề tài này, WebSocket được tận dụng để duy trì luồng truyền tải dữ liệu nặng, liên tục như truyền hình ảnh/video thô (video streaming) từ camera OV2640/OV3660 lên ứng dụng di động trong quá trình cấu hình khuôn mặt từ xa.

\textbf{MQTT (Message Queuing Telemetry Transport)} là giao thức truyền thông gọn nhẹ theo mô hình Xuất bản - Đăng ký (Publish-Subscribe) hoạt động trên nền tảng TCP/IP, được thiết kế chuyên dụng cho các thiết bị IoT tại biên có tài nguyên cực kỳ hạn chế và băng thông mạng không ổn định. Hệ thống sử dụng một MQTT Broker trung gian (như Mosquitto hoặc EMQX). Vi điều khiển ESP32-S3 sẽ \textit{Publish} các gói tin dữ liệu dạng chuỗi JSON cực nhẹ lên các topic như \texttt{lock/status} hoặc \texttt{lock/logs} để cập nhật trạng thái đóng/mở cửa hoặc gửi nhật ký sự kiện về Server ngay khi nhận diện xong. Ngược lại, khi người dùng bấm nút mở khóa trên ứng dụng di động, Server sẽ \textit{Publish} một lệnh điều khiển xuống topic \texttt{lock/control}. ESP32-S3 đã \textit{Subscribe} trước topic này sẽ ngay lập tức nhận được payload và kích hoạt chân GPIO điều khiển Relay mở chốt cửa điện từ với độ trễ chỉ vài mili-giây.

\subsubsection{Hệ quản trị cơ sở dữ liệu PostgreSQL}
Toàn bộ dữ liệu của hệ thống quản lý căn hộ và vận hành khóa thông minh được lưu trữ lâu dài tại Server bằng Hệ quản trị cơ sở dữ liệu quan hệ đối tượng (ORDBMS) PostgreSQL. Việc lựa chọn PostgreSQL làm nền tảng lưu trữ không chỉ đảm bảo tính toàn vẹn dữ liệu thông qua các ràng buộc tham chiếu chặt chẽ (Foreign Keys) và tuân thủ tuyệt đối các tính chất giao dịch ACID (Atomicity, Consistency, Isolation, Durability), mà còn mang lại nhiều lợi thế vượt trội về khả năng mở rộng và hiệu năng. PostgreSQL nổi tiếng với kiến trúc đa phiên bản đồng thời (MVCC), cho phép nhiều tiến trình đọc và ghi diễn ra song song mà không khóa lẫn nhau, tối ưu hóa cho các ứng dụng có lượng truy cập cao. Hơn nữa, PostgreSQL hỗ trợ mạnh mẽ các kiểu dữ liệu phức tạp như mảng (arrays), JSON/JSONB và các kiểu dữ liệu hình học không gian (PostGIS), tạo điều kiện thuận lợi để lưu trữ các cấu trúc dữ liệu đa dạng của hệ thống IoT mà không cần phải hy sinh các quy tắc chuẩn hóa của mô hình quan hệ.